\section{ADR-010: Nearest-rank percentiles for latency reporting}
\label{adr:010}

\subsection*{Status}

Accepted; in force at revision \texttt{\projectrevisionshort}.

\subsection*{Context}

The benchmark harness reports a 95th-percentile latency per pipeline phase.
Percentile definitions differ, and interpolating definitions can return a value
that was never observed.

\subsection*{Decision}

\texttt{benchmarking.percentile} implements the nearest-rank method:
$\rho = \lceil p n \rceil$ clamped to $[1, n]$, returning $y_{(\rho)}$ from the
sorted samples. The source gives the rationale directly: nearest-rank ``always
returns an actually-observed sample, which is easy to reason about for latency
SLOs and avoids fabricating values between measurements''.

\subsection*{Consequences}

\begin{itemize}
  \item Every reported percentile corresponds to a real measurement.
  \item The function is exact and has no floating-point interpolation, so it is
    trivially testable and reproducible.
  \item For small sample counts the estimate is coarse: with ten repetitions,
    $\operatorname{P}_{0.95}$ is the maximum. The report carries the sample
    count for every statistic, so the coarseness is visible rather than
    implied.
  \item Numbers are not directly comparable with tools that interpolate, such
    as NumPy's default \texttt{linear} method.
\end{itemize}

\subsection*{Alternatives}

Linear interpolation between order statistics (NumPy's default), or a
streaming estimator such as t-digest. The first fabricates unobserved values;
the second is disproportionate for a local benchmark harness.

\subsection*{Evidence}

\srcfile{src/feedback\_intelligence\_agent/benchmarking.py}
(\texttt{percentile} and its docstring; \texttt{summarize} reporting mean,
median, p95, min, max, and sample count);
\srcfile{tests/test\_benchmarking.py}.
